\section{Angiography}


\subsubsection*{What Is This Test?}
Angiography is an X-ray imaging technique that uses a radiopaque contrast dye injected into blood vessels to visualize blood flow and locate blockages.

\subsubsection*{Alternative Names}
Coronary angiography, Arteriogram, Vascular imaging

\subsubsection*{Principle}
Angiography is an X-ray imaging technique that uses a radiopaque contrast dye injected into blood vessels to visualize blood flow and locate blockages.

\subsubsection*{History}
Egas Moniz developed cerebral angiography in 1927 in Portugal. F. Mason Sones accidentally performed the first selective coronary angiogram in 1958.

\subsubsection*{Why This Test Is Ordered}
Ordered to visualize narrowing or blockages in the coronary, carotid, renal, or peripheral arteries, especially in patients with suspected coronary artery disease or stroke.

\subsubsection*{Is This a Primary or Secondary Test}
Primary diagnostic standard for defining arterial anatomy and stenosis.

\subsubsection*{Technical success and diagnostic accuracy}
Angiography may be technically completed without resolving the clinical problem. Report vessel visualization, treatment endpoint, complications, and follow-up as separate outcomes.

\subsubsection*{How This Test Is Performed}
Contrast dye is injected into the target artery via a catheter under continuous fluoroscopy (X-ray monitoring).

\subsubsection*{What Preparation Is Needed}
Preparation, renal-function review, and fasting depend on the planned angiographic procedure and local protocol. Review prior contrast reactions and medications.

\subsubsection*{Contraindications and Precautions}
Pregnancy is not an automatic contraindication; the examination should be justified and optimized, and alternatives considered when appropriate. For contrast-enhanced studies, review the indication, prior severe reactions, and renal function according to local protocol.

\subsubsection*{Risks}
Risks associated with X-ray imaging include:
\begin{itemize}
    \item \textbf{Radiation exposure}: Exposure to low levels of ionizing radiation, with cumulative lifetime dose optimization guided by the ALARA (As Low As Reasonably Achievable) principle.
    \item \textbf{Contrast reaction}: Allergic reactions can occur, including rare severe reactions. Kidney injury after iodinated contrast is uncommon and risk depends on baseline kidney function, illness, and the clinical context.
\end{itemize}


Drink plenty of water to eliminate contrast dye; monitor vital signs and renal function.

\subsubsection*{Understanding Your Results}
Identifies the location and severity of arterial narrowing (stenosis) graded as percentage occlusion.

\subsubsection*{Normal Results}
Smooth, patent blood vessels with no significant narrowing or stenosis.

\subsubsection*{Follow-up Tests}
Angioplasty, stenting, or vascular surgery consult.

\subsubsection*{References}
\begin{enumerate}
  \item MedlinePlus. Angiography. U.S. National Library of Medicine; 2025.
  \item Current Medical Diagnosis and Treatment. McGraw-Hill; 2025.
\end{enumerate}

\clearpage
